\documentclass[10pt,a4paper]{article}
\usepackage[margin=2.2cm]{geometry}
\usepackage{amsmath,amssymb}
\usepackage{booktabs}
\usepackage{longtable}
\usepackage[hidelinks]{hyperref}
\usepackage{parskip}
\usepackage{enumitem}
\setlist{nosep,leftmargin=1.4em}

\title{\textbf{Graph Sparsification and Community Detection}\\[2pt]
\large A chronological review of the literature, 1996--2026,\\
with an audit of what this project holds and what it does not}
\author{Compiled 2026-07-27}
\date{}

\begin{document}
\maketitle

\section*{Scope, and how to read the verification marks}

This covers one question and its immediate neighbourhood: \emph{does removing edges from a graph
before running community detection help, hurt, or leave it unchanged?} It is not a survey of
graph sparsification in general, nor of community detection in general. Works from those fields
appear only where the pipeline depends on them.

Every entry carries a verification mark, because the value of a review of this kind is entirely
in whether its statements can be trusted:

\begin{description}[leftmargin=2.6em,style=nextline]
\item[\textbf{[FULL]}] The paper was read end to end from the PDF, and any characterization of
what it claims can be relied on. Eight papers are in this class and the reading notes are in
\texttt{PAPER\_RESTRUCTURE/references/}.
\item[\textbf{[META]}] Bibliographic fields were read off a publisher record, DOI resolver,
arXiv, or the Semantic Scholar API. Safe to cite. Says nothing about whether the paper was read.
\item[\textbf{[2ND]}] Described from a fetched page or search summary, not read directly. Listed
as a lead. \emph{Not safe to characterize.}
\end{description}

Where a claim about a paper's content appears below without a \textbf{[FULL]} mark, it is a
statement about what the paper says it does, not a verified account of what it does.

\section{Prehistory: what sparsification guaranteed before anyone applied it to communities}

The theory arrived first and it never promised what the pipeline later assumed.

\textbf{Bencz\'ur and Karger (1996, journal 2015)}~\cite{benczur2015randomized} \textbf{[META]}
preserve every cut of a graph to within $1 \pm \epsilon$ by sampling edges with probability
proportional to connectivity and reweighting the survivors.

\textbf{Spielman and Teng (2004, journal 2014)}~\cite{spielman2004nearly, spielman2014nearly}
\textbf{[META]} and \textbf{Spielman and Srivastava (2008)}~\cite{spielman2008graph}
\textbf{[META]} preserve the Laplacian quadratic form by sampling according to effective
resistance~\cite{lovasz1993random, doyle1984random}, giving a spectral sparsifier with
$O(n \log n / \epsilon^2)$ edges.

\emph{What this does and does not license.} The guarantee is about the Laplacian spectrum.
Community structure is related to that spectrum through the eigenvectors that encode
cuts~\cite{von2007tutorial, shi2000normalized, newman2013spectral}, but it is not determined by
it. Every later claim that a spectral sparsifier ``preserves communities'' is an extrapolation
beyond the theorem, and the gap between the two is the space this literature occupies.

Two structural filters predate the pipeline and are still used as baselines: the \emph{disparity
filter} of \textbf{Serrano et al.\ (2009)}~\cite{serrano2009extracting} \textbf{[META]}, which
keeps edges improbable under a local null, and its noise-corrected successor from
\textbf{Coscia and Neffke (2017)}~\cite{coscia2017backboning} \textbf{[META]}.

\section{The founding claim: Satuluri, Parthasarathy and Ruan, SIGMOD 2011}

\textbf{[FULL]}~\cite{satuluri2011local}. This is the origin of the sparsify-then-detect
pipeline and everything downstream is a response to it. Read in full; notes in
\texttt{NOTES\_satuluri2011.md}.

The method, \emph{L-Spar}, ranks each edge by the Jaccard similarity of its endpoints'
neighbourhoods, computed via minhash signatures, and has each vertex retain its top
$\lceil d_v^{\,e}\rceil$ edges. The abstract claims speedups ``often in the range 10--50'' with
``little or no deterioration in the quality of the resulting clusters,'' and, decisively,
``for at least two of the four clustering algorithms, our sparsification consistently enables
higher clustering accuracies.''

\emph{Which two.} Transcribed from their Table 2 (p.~729), ground-truth F-score by algorithm:
Metis 4W/0L, Graclus 3W/0L, Metis+MQI 2W/2L, MLR-MCL 1W/4L. The accuracy claim is about the two
\textbf{fixed-$k$ cut partitioners}. The flow-based algorithm is the one L-Spar does not help.
This is the single most misquoted fact in the downstream literature.

\emph{What they got right, and what the field then lost.} Their \S4.3 states, verbatim, that
clusters obtained from a sparsified graph cannot be scored on that same graph, ``since that
would tell us nothing about how well the sparsified graph retained the cluster structure in the
original graph,'' and they report every conductance on the original graph. They fix the number
of communities once per network and give it to every arm, so granularity is matched by
construction. They charge sparsification time to the reported speedup. On all three counts the
1996--2011 origin paper is stricter than much of what followed it.

\emph{Their own scope conditions.} A synthetic sweep in their \S4.5 reports the method beginning
to outperform clustering on the unsparsified graph only from average degree $\approx 50$. Their
speedups are measured against clustering runs lasting hours, with an approximate minhash
similarity reported as two orders of magnitude cheaper than exact Jaccard. Their disclosed
limitation is that L-Spar cannot handle bipartite or triangle-free structure without a
bibliographic-coupling transform.

The two detectors the accuracy claim is about:
\textbf{Graclus}~\cite{dhillon2007weighted} \textbf{[META]}, multilevel weighted kernel
$k$-means for normalized cut, with no balance constraint; and
\textbf{Metis}~\cite{karypis1998metis} \textbf{[META]}, multilevel $k$-way partitioning with a
hard balance constraint. The flow-based one is
\textbf{MLR-MCL}~\cite{satuluri2009scalable} \textbf{[META]}. All three are distributed by their
authors and build on a current toolchain.

\section{The parallel question: does sampling preserve community structure at all?}

A separate lineage asks whether \emph{any} reduction preserves communities, independent of
whether it accelerates anything.

\textbf{Leskovec and Faloutsos (2006)}~\cite{leskovec2006sampling} \textbf{[META]} set the
graph-sampling problem. \textbf{Maiya and Berger-Wolf (2010)}~\cite{maiya2010sampling}
\textbf{[META]} sample explicitly to preserve community structure.

\textbf{Blagus et al.\ (2015)}~\cite{blagus2015sampling} \textbf{[META]} is the dissenting
result and deserves more attention than it gets: sampling can \emph{manufacture} the appearance
of community structure rather than preserve it. Any measured ``improvement'' after reduction has
to clear this before it can be called a gain.

\textbf{Peng, Kolda and Pinar (2014)}~\cite{peng2014accelerating} \textbf{[META]} reduce by
$k$-core before detecting, which removes vertices as well as edges and therefore scores its
partition on a smaller, denser vertex set.

\section{Adoption and extension, 2016--2021}

\textbf{Hamann et al.\ (2016)}~\cite{hamann2016structure} \textbf{[FULL]} is the systematic
comparison, and it released the implementations now distributed in standard graph libraries,
including Local Degree, Local Similarity and K-Neighbor. Two of its findings matter here and are
frequently ignored. First, verbatim: normalized mutual information ``reports higher similarity
values for a larger number of found communities \dots\ This makes it unsuitable for comparing
partitions on sparsified networks,'' with the adjusted Rand index recommended instead. Second,
that filtering ``does lead the algorithm into finding different community structures,'' and that
the preserved structure ``is not necessarily the same as the one the Louvain algorithm finds.''

\textbf{Sotiropoulos and Tsourakakis} contribute twice: a unifying framework for
spectrum-preserving sparsification and coarsening at NeurIPS
2019~\cite{sotiropoulos2019unifying} \textbf{[META]}, and triangle-aware spectral sparsifiers at
KDD 2021~\cite{sotiropoulos2021triangle} \textbf{[FULL]}. The 2021 paper is worth reading for
its framing alone: it treats sparsification as an established and effective route to community
detection whose \emph{success requires explanation}, opening ``how can we explain its empirical
success?''\ By 2021 the pipeline is a premise, not a proposition.

\textbf{Satuluri et al.\ (2020)}~\cite{satuluri2020simclusters} \textbf{[META]} put
sparsification before community detection in production at Twitter.
\textbf{Wu and Chen (2020)}~\cite{wu2020gsgan} \textbf{[META]} train a GAN to produce the
subgraph.

\textbf{Laeuchli (2020)}~\cite{laeuchli2020fast} \textbf{[FULL]} is the one analysis that
identifies a regime where the approach provably pays: dense planted-partition graphs with
spectral solvers, where detection cost is superlinear in the number of edges and the planted
structure survives heavy deletion. It is the honest boundary of the pipeline's scope, and it is
not the regime the pipeline is usually applied to.

\section{The GNN turn, and what it cost this question}

\textbf{Liu et al.\ (2023)}~\cite{liu2023dspar} \textbf{[FULL]} introduce \emph{DSpar},
approximating the effective resistance of $(u,v)$ by $1/d_u + 1/d_v$ and thereby obtaining a
spectral sparsifier at the cost of reading the degree sequence. It is designed for GNN training,
not community detection, and its adoption for the latter is downstream extrapolation.

\textbf{Hashemi et al.\ (IJCAI 2024)}~\cite{hashemi2024survey} \textbf{[FULL]} survey graph
reduction across sparsification, coarsening and condensation. The most useful fact about this
survey is negative and was established by text search: it contains four occurrences of ``graph
classification,'' one of ``node classification,'' one of ``clustering,'' and \textbf{zero} of
``community.'' The modern reduction field has reoriented entirely around GNN prediction and has
dropped community detection as an evaluation target. The question was asked in 2011--2016,
answered optimistically, and then abandoned rather than settled.

\section{The benchmark era: Chen et al., PVLDB 2024}

\textbf{[FULL]}~\cite{chen2024demystifying}, notes in \texttt{NOTES\_chen\_and\_fastcd.md}.
Twelve sparsifiers, sixteen graph properties, fourteen networks. Conclusions: no single
sparsifier preserves all properties; the choice should follow the downstream task; pruning
fragments graphs; and fidelity to the full graph's partition degrades monotonically with the
prune rate.

Two things about this paper matter disproportionately.

\emph{What it owns.} Fragmentation-with-pruning at benchmark scale, and monotone fidelity
degradation, are established facts of this literature and belong to Chen et al.

\emph{What it does not measure.} It never computes a quality objective for community detection.
Its clustering metric is fidelity to the full graph's partition, not whether the sparsified
partition is \emph{better or worse}. It also never times the downstream task. So the field's two
central premises, that sparsification preserves quality and delivers speed, are both untested in
the field's most rigorous benchmark. Further, its released code computes modularity for the
sparsified condition \emph{on the sparsified graph}
(\texttt{metrics\_nk.py:414}), which is the clearest available evidence that the
scoring hazard is present in the tooling of the most careful work in the area.

\section{Redundancy: why any of this is possible}

A separate and largely disconnected strand explains why graphs can lose most of their edges at
all.

The \emph{metric backbone} of a weighted graph is the union of its all-pairs shortest paths; an
edge outside it is redundant in the exact sense that no shortest path uses it.
\textbf{Correia, Barrat and Rocha (2023)}~\cite{correia2023contact} \textbf{[META]} report the
backbone to be a small fraction of the edges of real contact networks while community structure
is retained. \textbf{Dreveton, Chucri, Grossglauser and Thiran (NeurIPS 2024)}
\cite{dreveton2024metric} \textbf{[META, abstract read verbatim]} prove that deleting every edge
outside the backbone preserves community structure on a broad class of weighted random graphs
with communities.

\textbf{He, Drineas and Khanna (2025)}~\cite{he2025structure} \textbf{[META, abstract read
verbatim]} prove that under strong clusterability, characterized by a structure ratio
$\Upsilon(k) = \lambda_{k+1}/\rho_G(k)$, \emph{uniform} edge sampling preserves the spectral
subspace used for clustering, bypassing importance sampling entirely. If this holds empirically
it is a significant negative for the whole similarity-based line: the expensive selection rule
would be unnecessary where the structure is strong, and where it is weak the guarantee lapses.
The paper reports no empirical evaluation on community detection.

\section{The current wave, 2024--2026}

The pipeline is not historical. Five works in three years, all proposing or applying it.

\textbf{Socievole and Pizzuti} publish twice: ASONAM 2024~\cite{socievole2024community}
\textbf{[META]} and \emph{Soft Computing} 30:2109--2133, 2026~\cite{socievole2026effective}
\textbf{[FULL]}, notes in \texttt{NOTES\_socievole\_pizzuti\_2026.md}. The 2026 method weights a
graph by effective resistance, applies a kernel, sparsifies by weight thresholding, and runs a
genetic algorithm maximizing weighted modularity on the result. It is worth reading for three
reasons. Its evaluation measure is NMI against ground truth rather than a self-scored objective,
so it does not commit the scoring error. It compares against unsparsified Louvain and Infomap
throughout. And its own tables show those unsparsified baselines matching or beating every
sparsified variant it tests: at mixing $0.3$ on 1000-node LFR, Louvain reaches NMI $1$ against
the sparsified method's $0.601$, and its \S6.2 states that on the densest network ``OmeGANet
achieves only 0.3522 as NMI value, while Louvain and Infomap are able to perfectly match the
ground-truth.'' The gains it reports are of kernels and the genetic algorithm over its own
no-kernel baseline, not of sparsification over no sparsification.

\textbf{Setiadi, Yaakub and Abu Bakar (2025)}~\cite{setiadi2025community} \textbf{[FULL]},
$k$-core community-preserving sparsification, notes in \texttt{NOTES\_kcore\_ijain2026.md}.
Reproduced and falsified in this project. Its reported modularity of $0.63$ on Zachary's karate
club exceeds $0.4198$, the maximum attainable on that graph over all partitions, which
identifies the scored graph as the sparsified one.

\textbf{Pari, Bhandari and Raha (2026)}~\cite{pari2026effective} \textbf{[FULL]},
effective-resistance sparsification and community detection, notes in
\texttt{NOTES\_effres\_2026.md}. Reports modularity $0.87$ on the American College Football
network, whose maximum over all partitions is $0.6046$. Same diagnosis.

\textbf{Gottesb\"{u}ren et al.\ (ESA 2025)}~\cite{gottesburen2025linear} \textbf{[META]} come
from graph partitioning rather than community detection and observe the cost point directly:
sparsification overhead can cancel the speedup it is meant to produce.

\section{What the corpus contains, and what it does not}

\subsection*{Holdings}

\begin{center}
\begin{tabular}{@{}lr@{}}
\toprule
Papers read end to end, with written notes & 8 \\
PDFs archived locally & 8 \\
Lines of verified reading notes & 2{,}664 \\
Bibliography entries, all fields from fetched records & 72 \\
On-subject works spanning 2011--2026 & 14 \\
\bottomrule
\end{tabular}
\end{center}

The eight read in full are Satuluri et al.\ 2011 (plus the author's thesis), Hamann et al.\ 2016,
Laeuchli 2020, Sotiropoulos and Tsourakakis 2021, Liu et al.\ 2023, Chen et al.\ 2024, Hashemi et
al.\ 2024, Setiadi et al.\ 2025, Pari et al.\ 2026, and Socievole and Pizzuti 2026. Four of these
were not merely read but \emph{reproduced}, and two of the reproductions falsified the paper's
central claim.

\subsection*{Genuine gaps, in priority order}

\begin{enumerate}
\item \textbf{He, Drineas and Khanna (2025) is unread.} \textbf{[META]} only. If uniform sampling
provably suffices under strong clusterability, it bears directly on whether any similarity-based
selection rule is necessary. This is the highest-value unread paper in the area.
\item \textbf{Sotiropoulos and Tsourakakis (NeurIPS 2019) is unread.} \textbf{[META]} only. The
unifying spectrum-preserving framework is prior art for several framing choices.
\item \textbf{Dreveton et al.\ (2024): the proof is verified, the experiments are not.} The
abstract was read verbatim; the evaluation protocol was described only through a page summary.
Whether they control granularity under Leiden is unknown.
\item \textbf{Socievole and Pizzuti (ASONAM 2024) is unread.} The 2026 paper is its extension.
\item \textbf{The 2012--2015 window is thin} on this specific question. Repeated searches return
the same core corpus, which is consistent with the Hashemi survey's finding that the question was
largely dormant, but absence of search results is weak evidence.
\item \textbf{Coarsening and condensation applied to community detection} is covered only through
the Hashemi survey and Peng et al. If the pipeline has a modern descendant it is here.
\end{enumerate}

\subsection*{Assessment}

For the purpose of grounding a paper on this question, the corpus is close to complete on the
\emph{lineage} and thin on the \emph{theory adjacent to it}. Every work that proposes or applies
sparsify-then-detect and that a reviewer is likely to name has been located, and the five that
matter most have been read in full and in four cases reproduced. What is missing is not more
instances of the pipeline; searching for those has reached saturation, with the same dozen works
returning across many differently-phrased queries. What is missing is the recent theory that
constrains it, items 1 to 3 above, and of those only He et al.\ could change how the question is
posed.

\emph{Recommended additional work: read three papers.} He et al.\ 2025, Sotiropoulos and
Tsourakakis 2019, and the experimental section of Dreveton et al.\ 2024. That is the whole of
what a further literature effort should consist of. Beyond those three, further searching is
unlikely to return anything the corpus does not already hold.

\bibliographystyle{plain}
\bibliography{refs}

\end{document}
